\documentclass[10pt,a4paper]{amsart}
\usepackage[margin=19mm]{geometry}
\usepackage{amsmath,amssymb,mathtools}
\usepackage[scr=rsfs,cal=euler]{mathalfa}
\usepackage{xfrac}
\usepackage[unicode,hidelinks,bookmarks=false]{hyperref}
\newcommand{\ie}{i.e.\ }
\newcommand{\cf}{cf.\ }
\newcommand{\resp}{respectively\ }
\newcommand{\Subsection}{\subsection}
\providecommand{\nobreakdash}{\nobreak\hbox{-}}
\setlength{\emergencystretch}{2em}
\multlinegap0pt
\usepackage{latexsym}
\allowdisplaybreaks
\usepackage{bbm}
\usepackage{mathtools}
\usepackage{ulem}
\usepackage{tikz-cd}
\usepackage{mathrsfs}
\usepackage{todonotes}
\usepackage{enumitem}
\usepackage{verbatim}
\usepackage{amsaddr}
\usepackage{dsfont}
\newtheorem{thm}{Theorem}[section]
\newtheorem*{thm*}{Theorem}
\newtheorem{Con}[thm]{Conjecture}
\newtheorem{cor}{Corollary}[section]%Corollary单独排序，不和其他混淆
\newtheorem{prop}{Proposition}
\newtheorem{lem}[thm]{Lemma}
\newtheorem{pro}[thm]{Proposition}
\newtheorem{defin}{Definition}
\newtheorem{exa}[thm]{Example}
\newtheorem{rem}{Remark}[section]%Remark单独排序，不和其他混淆
\newtheorem*{xrem}{Remark}%Remark无编号
\numberwithin{equation}{section}
\newcommand{\bg}{\big}
\newcommand{\bgl}{\big(}
\newcommand{\bgr}{\big)}
\newcommand{\bgg}{\bigg}
\newcommand{\bggl}{\bigg(}
\newcommand{\bggr}{\bigg)}
\newcommand{\Bg}{\Big}
\newcommand{\Bgl}{\Big(}
\newcommand{\Bgr}{\Big)}
\newcommand{\Bgg}{\Bigg}
\newcommand{\Bggl}{\Bigg(}
\newcommand{\Bggr}{\Bigg)}
\newcommand{\lo}{\log_2}
\newcommand{\lt}{\log_3}
\newcommand{\inv}{^{-1}}
\newcommand{\ph}{\Phi}
\newcommand{\mbc}{\mathbb{C}}
\newcommand{\mbd}{\mathbb{D}}
\newcommand{\mbe}{\mathbb{E}}
\newcommand{\mbt}{\mathbb{T}}
\newcommand{\mbz}{\mathbb{Z}}
\newcommand{\mbr}{\mathbb{R}}
\newcommand{\mbs}{\mathbb{S}}
\newcommand{\mbf}{\mathbb{F}}
\newcommand{\mbh}{\mathbb{H}}
\newcommand{\mbl}{\mathbb{L}}
\newcommand{\mbn}{\mathbb{N}}
\newcommand{\mbp}{\mathbb{P}}
\newcommand{\mbq}{\mathbb{Q}}
\newcommand{\mbb}{\mathbb{B}}
\newcommand{\mbm}{\mathbb{M}}
\newcommand{\mbbx}{\mathbbm{x}}
\newcommand{\mbby}{\mathbbm{y}}
\newcommand{\mch}{\mathcal{H}}
\newcommand{\mca}{\mathcal{A}}
\newcommand{\mcb}{\mathcal{B}}
\newcommand{\mcc}{\mathcal{C}}
\newcommand{\mcd}{\mathcal{D}}
\newcommand{\mce}{\mathcal{E}}
\newcommand{\mcf}{\mathcal{F}}
\newcommand{\mcg}{\mathcal{G}}
\newcommand{\mcj}{\mathcal{J}}
\newcommand{\mcpg}{\mathcal{g}}
\newcommand{\mcph}{\mathcal{h}}
\newcommand{\mcn}{\mathcal{N}}
\newcommand{\mcp}{\mathcal{P}}
\newcommand{\mck}{\mathcal{K}}
\newcommand{\mcl}{\mathcal{L}}
\newcommand{\mco}{\mathcal{O}}
\newcommand{\mcq}{\mathcal{Q}}
\newcommand{\mcm}{\mathcal{M}}
\newcommand{\mcr}{\mathcal{R}}
\newcommand{\mcs}{\mathcal{S}}
\newcommand{\mfa}{\mathfrak{A}}
\newcommand{\mfpa}{\mathfrak{a}}
\newcommand{\mfb}{\mathfrak{B}}
\newcommand{\mfg}{\mathfrak{G}}
\newcommand{\mfr}{\mathfrak{R}}
\newcommand{\mfpg}{\mathfrak{g}}
\newcommand{\mfph}{\mathfrak{h}}
\newcommand{\mfz}{\mathfrak{z}}
\newcommand{\mse}{\mathscr{E}}
\newcommand{\msk}{\mathscr{K}}
\newcommand{\mskt}{\mathscr{K}_\tau}
\newcommand{\msu}{\mathscr{U}}
\newcommand{\mst}{\mathscr{T}}
\newcommand{\mckt}{\mathcal{K}_\tau}
\newcommand{\mmd}{\mathrm{d}}
\newcommand{\mme}{\mathrm{e}}
\newcommand{\mmi}{\mathrm{i}}
\newcommand{\mit}{\mathrm{i}t}
\newcommand{\re}{\operatorname{Re}}
\newcommand{\im}{\operatorname{Im}}
\newcommand{\me}{\operatorname{meas}}
\newcommand{\whp}{\widehat{\Phi}}
\newcommand{\ol}{\overline}
\hypersetup{colorlinks=true,linkcolor=blue,anchorcolor=blue,citecolor=blue}
\usepackage{color}
\DeclareMathOperator\dif{d\!}
\newcommand{\newabstract}[1]{%
	\par\bigskip
	\csname otherlanguage*\endcsname{#1}%
	\csname captions#1\endcsname
	\item[\hskip\labelsep\scshape\abstractname.]
}
\begin{document}
\title[Extreme values of derivatives of Dirichlet L-functions]{Extreme values of derivatives of Dirichlet $L$-functions}
\author{Xinghua Cui, Zhifeng Peng, Yutong Song, Shengbo Zhao}
\date{}
\maketitle
\noindent\emph{Source and licence.} Extreme values of derivatives of Dirichlet $L$-functions. Version arXiv:2607.04084v1 (CC BY 4.0). This material is distributed under the Creative Commons Attribution 4.0 International licence, http://creativecommons.org/licenses/by/4.0/, as recorded in the arXiv OAI licence metadata for this exact version and shown on the abstract page https://arxiv.org/abs/2607.04084v1. Full rights notice: licenses/arxiv-2607.04084-CC-BY-4.0.txt.\par\smallskip
\noindent\textit{Editorial scope.} Counted unit: the original \texttt{lem} environment \texttt{lem1} at source lines 233--239 together with its complete proof at lines 240--269; the delivered document keeps source lines 229--270 so that every referenced label and every citation resolves inside it. Native editable source is the paper's own concatenated TeX; the AMS wrapper is editorial formatting only. No publisher class, upstream script or unrelated artwork is redistributed.
\section{Preliminaries}
\label{sec-pre}
In this section, we carry out some preliminary work. We begin with the following asymptotic formula, which shows that derivatives of Dirichlet \(L\)-functions can be effectively approximated by a truncated Dirichlet series.


\begin{lem}
\label{lem1}
    Fix \(\ell \in \mbn_+\) and \(\sigma \in (1/2,1)\). When \(q\) is sufficiently large, for any non-principal character \(\chi\pmod {q}\), we have
    \[
    (-1)^\ell  L^{(\ell)}(\sigma,\chi) = \sum_{n \le q}\frac{(\log n )^\ell \chi(n)}{n^\sigma} + O_{\sigma,\ell} \bg( q^{\frac{1}{2}-\sigma}(\log q)^{\ell+1}\bg).
    \]
\end{lem}

\begin{proof}
    Trivially, we have
\begin{align}
    \label{lem1-dirichlet}
    (-1)^\ell  L^{(\ell)}(\sigma,\chi) = \sum_{n\in \mbn} \frac{(\log n )^\ell \chi(n)}{n^\sigma} = \Bg(\sum_{n \le q} +\sum_{n > q}\Bg) \frac{(\log n )^\ell \chi(n)}{n^\sigma}.
\end{align}
Set
\[
A_\chi (t) = \sum_{n \le t} \chi(n) ~\text{and}~ f(t) = \frac{(\log t)^\ell}{t^\sigma}.
\]
The P\'olya-Vinogradov inequality gives, uniformly for \(t \ge 1\), 
\begin{align}
    \label{lem1-pv}
    |A_\chi (t)| \ll \sqrt{q}\log q.
\end{align}
Hence, an application of partial summation shows that the tail sum on the right-hand side of \eqref{lem1-dirichlet} satisfies
\begin{align}
    \label{lem1-tail}
    \sum_{n > q}\frac{(\log n )^\ell \chi(n)}{n^\sigma} = -A_\chi (q) f(q) + \int_{q}^\infty A_\chi(t) f^\prime (t) \mmd t. 
\end{align}
Noting that
\[
\bg|f^\prime(t)\bg| \ll_{\sigma,\ell} \frac{(\log t)^\ell}{t^{\sigma+1}}.
\]
Combining this with \eqref{lem1-pv} and \eqref{lem1-tail}, we obtain
\[
\sum_{n > q}\frac{(\log n )^\ell \chi(n)}{n^\sigma} \ll_{\sigma,\ell} q^{\frac{1}{2}-\sigma}(\log q)^{\ell+1}.
\]
Substituting this into \eqref{lem1-dirichlet} completes the proof of Lemma \ref{lem1}.
\end{proof}


\end{document}
